\documentclass[11pt]{article}
\usepackage{amsmath,amssymb,fullpage}

\title{Why F2 Pruning Fails Above $\sim$2.5\% Error at Rates $R>0.5$}
\author{Analysis Note}
\date{\today}

\begin{document}
\maketitle

\section{Setup}

The pruning decoder works in two stages:
\begin{enumerate}
  \item \textbf{Column selection:} From $n$ columns, pick those whose row-support has
        high intersection with the syndrome support $\mathcal{S}=\{r : s_r\neq0\}$.
  \item \textbf{Solve:} Build a dense $m\times k$ system ($k\approx m-50$ overdetermined)
        and solve via Gauss--Jordan elimination.
\end{enumerate}

Success requires both stages to work: (i) the true error columns must be selected,
and (ii) the resulting linear system must have full column rank.

\section{Syndrome statistics}

Let $H$ be an $m\times n$ binary parity-check matrix with column weight $c$ and
row weight $r=c/R$.  Transmit the all-zero codeword over a BSC with bit-error
rate $\beta$.  For a particular check node (row) $i$, the syndrome bit is

\[
s_i = \bigoplus_{j\;:\;H_{ij}=1} e_j .
\]

For a fixed row $i$ of weight $r$, the $r$ incident bits are each 1 with
probability $\beta$ independently.  Therefore

\[
P(s_i = 1) = \frac{1 - (1-2\beta)^r}{2}.
\]

For small $\beta$,
\[
P(s_i = 1) \approx \frac{1 - e^{-2r\beta}}{2}.
\]

\section{Column selection: true positives vs false positives}

Consider a column $j$ that is \emph{not} in error.  Its $c$ incident rows each
have syndrome $1$ with probability $p_\text{odd}=P(s_i=1)$.  The number of its
rows in the syndrome support is $\mathrm{Binomial}(c, p_\text{odd})$.  With
threshold $v$, the probability of a false positive is

\[
p_f = P\bigl(\mathrm{Bin}(c, p_\text{odd}) > v\bigr) .
\]

With $c=6$ and $v=5$, this requires $6$ out of $6$ rows to have syndrome $1$:
\[
p_f = p_\text{odd}^6 .
\]

The expected number of false-positive columns is $n\cdot p_f$.
At $\beta=2.5\%$, $R=0.51$, $r=c/R\approx 12$:

\[
p_\text{odd} \approx \frac{1-e^{-2\cdot12\cdot0.025}}{2}
           = \frac{1-e^{-0.6}}{2} \approx 0.226,
\qquad
p_f \approx 0.226^6 \approx 1.3\times10^{-4}.
\]

Expected false positives: $2000\times 1.3\times10^{-4} \approx 0.26$ —
well under 1.  So false positives are not the bottleneck.

\section{The real bottleneck: rank deficiency from cancellation}

The critical issue for binary codes is \textbf{syndrome cancellation}.
At a single check node of degree $r$, the probability that it has an
\emph{even} number of incident errors (so $s_i=0$ despite errors) is

\[
P(\text{even}\mid \text{errors present})
= \frac{1 - (1-2\beta)^r}{2(1-(1-\beta)^r)} .
\]

For $\beta=0.025$, $r=12$: $P(\text{even}\mid\text{errors present}) \approx 0.39$.
Nearly $40\%$ of the time, an affected check node reports a zero syndrome.

This has two consequences:

\begin{enumerate}
  \item \textbf{Intersection shrinkage:} A column in error expects only
        $c\cdot p_\text{odd} \approx 6 \times 0.226 \approx 1.4$ of its $6$
        rows to have non-zero syndrome.  With $v=5$ we need $\ge 6$, so
        \emph{very few error columns clear the threshold}.  The pruning
        stage often selects too few true-error columns, forcing the
        candidate list into mostly irrelevant columns.

  \item \textbf{Effective rank loss:} The $m\times k$ dense system is
        nominally overdetermined by $m-k = 50$ rows.  But cancellation
        means many rows carry little or no information about the true
        errors.  The \emph{effective rank} of the system is lower than
        $k$, so Gauss--Jordan either finds a singular pivot or an
        inconsistent row.
\end{enumerate}

\section{Why larger fields don't suffer this}

For GF(256) each non-zero entry $H_{ij}\in\text{GF}(256)^\times$ is uniformly
random.  The syndrome at row $i$ is

\[
s_i = \sum_{j} H_{ij}\, e_j \quad (\text{in GF(256)}),
\]

so cancellation of two errors at the same check requires
$H_{i j_1}e_{j_1} + H_{i j_2}e_{j_2} = 0$, i.e.\
$e_{j_2} = e_{j_1}\cdot H_{i j_1}/H_{i j_2}$.  With uniformly random
field values, this is a probability-$1/255$ event — essentially never.
Hence \emph{no} syndrome cancellation, the intersection test is reliable,
and the dense system has full rank.

\section{Quantitative condition for F2 viability}

The pruning decoder succeeds reliably when:

\[
c \cdot p_\text{odd} > v \qquad\text{(error columns are selected)}
\]

and

\[
m - (m - 50) = 50 \gg \text{rank loss due to cancellation}.
\]

The first condition with $c=6$, $v=5$ requires $p_\text{odd} > 5/6 \approx 0.83$,
which would need

\[
\frac{1-e^{-2r\beta}}{2} > 0.83 \;\Longrightarrow\; e^{-2r\beta} < -0.66
\]

which is impossible.  In practice the decoder relies on the \emph{candidate
rotation} strategy: when sure columns are scarce, it picks the $k$ best
candidates and tries random subsets.  With $c=6$ the intersection distribution
is very coarse (only $7$ possible values $0,\dots,6$).  At $\beta=2.5\%$ the
error columns barely reach intersection $6$ (needed for $v=5$), so they fall
into the candidate pool where random rotation finds them some of the time.
At $\beta=3\%$ the distribution has shifted further: too few error columns
have $6$ intersections, and too many non-error columns do.

The empirical threshold of $\beta\approx 2.7$--$2.8\%$ for $R=0.51$ matches
the point where $p_\text{odd} \approx 0.24$, giving
$P(\mathrm{Bin}(6,0.24)=6) \approx 2\times10^{-4}$ — the intersection
distribution for error and non-error columns becomes nearly indistinguishable.

\section{Summary}

\begin{itemize}
  \item \textbf{Syndrome cancellation} (even number of errors at a check)
        is the fundamental limitation for F2 pruning.
  \item With $c=6$, the intersection test has only $7$ possible values,
        limiting discrimination.
  \item At $R=0.51$ ($r\approx12$), the maximum viable error rate is
        $\beta\approx 2.7$--$2.8\%$.
  \item Below $R=0.5$ there are more rows → more overdetermination →
        more tolerance for cancellation — hence augmentation works at
        $3$--$4\%$.
  \item GF(16) and GF(256) avoid cancellation entirely because random
        field values make weighted-sum collisions a $1/(q-1)$ event.
\end{itemize}

\end{document}
